\documentclass{article}
\usepackage[T1]{fontenc}
\usepackage{lmodern}
\usepackage{graphicx}
\usepackage{amsmath,amssymb}
\usepackage[a4paper,margin=2cm]{geometry}
\usepackage[absolute,overlay]{textpos}
\setlength{\TPHorizModule}{1mm}
\setlength{\TPVertModule}{1mm}
\usepackage[indonesian]{babel}
\usepackage{hyperref}
\setlength{\emergencystretch}{2em}
% unit: Halaman 21

\begin{document}

% === Halaman 21 (llm) ===

\begin{itemize}
  \item Untuk menggunakan OTP, pengirim dan penerima harus menyepakati kunci yang digunakan, kunci harus sepanjang pesan, dan hanya boleh digunakan sekali.
  \item Hal ini sulit dilakukan secara praktek.
  \item \textit{OTP} hanya dapat digunakan jika tersedia saluran komunikasi kedua yang sangat aman untuk mengirim kunci.
  \item Saluran kedua ini tidak boleh sama dengan saluran untuk mengirim pesan.
  \item Saluran kedua ini umumnya lambat dan mahal (misalnya lewat jalur darat, memakai kurir terpercaya dan tidak bisa dikenali).
  \item Atau, kunci dikirim pada saluran yang sama dengan saluran pengiriman pesan tetapi pada waktu yang berbeda dengan pengiriman pesan.
\end{itemize}

\end{document}
